\section{Pole bounds for bar faces and ordered collisions}
\label{op:sec:ordered-poles}

The multiplication faces of an internal bar compose adjacent inputs.
For coefficient operations with poles, this requires a filtration
preserved by the actual multiplication.  A collision expansion has its
own bound.  Already three
points give
\begin{equation}\label{op:eq:three-simple-poles}
 [t^{-1}]\frac{1}{t(x+t)x}=\frac{1}{x^2},
 \qquad z_1=z_2+t,\quad x=z_2-z_3.
\end{equation}
Thus a residue can merge three simple poles into a double pole.
Multiplication in an associative bar has a different, additive bound.
We first describe that bound, then calculate the collision coefficients.
All coordinates and series in this section are algebraic or formal.

\subsection{The filtration preserved by multiplication faces}

Let \(I\) be a finite labelled set, with an auxiliary total order used
to write each unordered difference once.  Put
\[
 R_I=k[z_i:i\in I],\qquad
 L_I=R_I[(z_i-z_j)^{-1}:i<j].
\]
For a tuple \(N=(N_{ij})_{i<j}\) of nonnegative integers define
\[
 D_I(N)=\left(\prod_{i<j}(z_i-z_j)^{-N_{ij}}\right)R_I.
\]
This is a vector subspace of \(L_I\), rather than a subalgebra.
It consists of rational functions with pole order at most \(N_{ij}\)
along each diagonal.  Distinct unordered differences are nonassociated
prime elements of the polynomial ring.  The pole bounds therefore have
their usual valuation meaning, independently of the chosen ordering.

Choose distinct \(\lambda_i\in k\).  Evaluation at \(z_i=\lambda_i\)
defines an augmentation \(\varepsilon:L_I\to k\).  Let
\(\bar L_I=\ker\varepsilon\) and
\(\bar D_I(N)=D_I(N)\cap\bar L_I\).  The normalized bar has
\(B_r=\bar L_I^{\otimes r}\), with homological differential
\begin{equation}\label{op:eq:pole-bar-differential}
 b[f_1|\cdots|f_r]
 =\sum_{i=1}^{r-1}(-1)^i
 [f_1|\cdots|f_if_{i+1}|\cdots|f_r].
\end{equation}
Here all algebra elements are in degree zero and have even parity.
Associativity makes \(b^2=0\), and the outer faces vanish on the
augmentation ideal.

\begin{proposition}[Total pole bounds in the normalized bar]
\label{op:prop:bar-pole-budget}
For a nonnegative tuple \(P\), set
\[
 F_PB_r=
 \sum_{N^{(1)}+\cdots+N^{(r)}\le P}
 \bar D_I(N^{(1)})\otimes\cdots\otimes\bar D_I(N^{(r)}),
 \qquad F_PB_0=k.
\]
The inequalities are componentwise.  These spaces form an exhaustive
increasing filtration by subcomplexes.  The multiplication bound
\[
 D_I(N)D_I(M)\subseteq D_I(N+M)
\]
is sharp on every diagonal.  A bound imposed separately on every bar
entry need not define a subcomplex.
\end{proposition}

\begin{proof}
Write the two factors with the displayed denominators.  Multiplication
adds their exponent tuples, and a product of augmentation-ideal
elements remains in that ideal.  In each term of
\eqref{op:eq:pole-bar-differential}, replace
\(N^{(i)},N^{(i+1)}\) by \(N^{(i)}+N^{(i+1)}\).
The total tuple does not change.  This proves preservation of \(F_P\).
Every tensor is a finite sum of pure tensors whose rational entries
have finite pole bounds, proving exhaustivity.

For sharpness it suffices to use
\(L=k[d,d^{-1}]\), with \(\varepsilon(d)=1\).
For any \(N\ge1\), the element \(u_N=d^{-N}-1\) belongs to
\(\ker\varepsilon\) and has exact pole order \(N\) at \(d=0\).
The only inner face of \([u_N|u_N]\) is
\(-[u_N^2]\), and
\[
 u_N^2=d^{-2N}-2d^{-N}+1
\]
has exact pole order \(2N\).  More generally \(u_Nu_M\) has exact
order \(N+M\).  Thus a uniform positive bound on individual entries
fails under an actual bar face.
\end{proof}

\begin{remark}[Linear filtrations with several divisors]
\label{op:rem:linear-pole-filtration}
For a coordinate divisor \(E\), let \(v_E\) be its discrete valuation,
normalized to take value one on a local equation.  Put
\(p_E(f)=\max(0,-v_E(f))\) for \(f\ne0\), and \(p_E(0)=0\).
The valuation identity gives
\[
 p_E(fg)\le p_E(f)+p_E(g).
\]
Equality holds when both valuations are nonpositive.  A zero of one
factor can reduce the pole of their product.

A single integer bound for several divisors gives a linear bar
filtration by taking
\[
 F_m^{\mathrm{tot}}B_r=\sum_{\sum_{i<j}P_{ij}\le m}F_PB_r,
 \qquad m\ge0.
\]
Each summand is a subcomplex by
Proposition~\ref{op:prop:bar-pole-budget}, so their sum is a subcomplex.
The linear span is necessary.  In \(k[x^{\pm1},y^{\pm1}]\),
augmented by evaluation at \((1,1)\), let
\[
 u=x^{-1}-1,\qquad v=y^{-1}-1.
\]
Each has total pole order one along the coordinate divisors.
Their sum has order one along each divisor, hence total order two.
The respective leading coefficients are one, so this remains true
in characteristic two.  Thus the literal set of functions with
bounded sum of divisor pole orders need not be a vector subspace.
The filtration above takes sums of the separately bounded spaces.
\end{remark}

The two-sided simplicial bar of
Proposition~\ref{prop:internal-bar} uses multiplication at fixed Ran
arity.  Its geometric definition supplies no rational-coordinate
presentation or pole lattice.  Proposition~\ref{op:prop:bar-pole-budget}
applies to such a presentation when multiplication respects the stated
filtration.  A collision of configuration labels is a separate map.

\subsection{An exact bound for a collision coefficient}

Relabel the colliding pair as \(1,2\), retaining label \(2\).
For \(j\ge3\) put
\[
 A=N_{12},\qquad B_j=N_{1j},\qquad C_j=N_{2j},\qquad
 x_j=z_2-z_j.
\]
Substitute \(z_1=z_2+t\) and expand at \(t=0\), with every
\(x_j\) invertible before completion.  This defines a ring map
\[
 \operatorname{Exp}_{12}:L_I\longrightarrow
 L_{I\setminus\{1\}}[[t]][t^{-1}].
\]
Indeed, \(z_1-z_2\) maps to \(t\), and
\(z_1-z_j\) maps to the unit \(x_j+t\) in the power-series ring.
Let \(C_{12}^q=[t^q]\operatorname{Exp}_{12}\).  This is a linear
coefficient map, not a ring map.  Its residue convention is
\(\operatorname{Res}_{t=0}f(t)\,dt=C_{12}^{-1}(f)\).

\begin{theorem}[Sharp collision denominator bound]
\label{op:thm:collision-poles}
Let \(f\in D_I(N)\).  Then \(C_{12}^q(f)=0\) for \(q<-A\).
For \(r=q+A\ge0\), put \(S=\{j\ge3:B_j>0\}\).  The
coefficient \(C_{12}^q(f)\) belongs to the finite sum
\begin{equation}\label{op:eq:collision-budget}
 \sum_{\substack{h\in\mathbb N^S\\\sum_{j\in S}h_j\le r}}
 \left(\prod_{3\le i<j}(z_i-z_j)^{-N_{ij}}\right)
 \left(\prod_{j\ge3}x_j^{-(B_j+C_j+h_j)}\right)
 k[z_2,\ldots,z_{|I|}],
\end{equation}
where \(h_j=0\) outside \(S\).  In particular, the pole order
against the surviving point \(j\) is at most
\begin{equation}\label{op:eq:pair-collision-bound}
 M_{2j}=
 \begin{cases}
 B_j+C_j+r,&B_j>0,\\
 C_j,&B_j=0.
 \end{cases}
\end{equation}
All other pair bounds remain \(N_{ij}\).  Each bound in
\eqref{op:eq:pair-collision-bound} is the least bound valid for every
polynomial numerator with the specified denominator tuple.
\end{theorem}

\begin{proof}
Write
\[
 f=\frac{P(z)}{\prod_{i<j}(z_i-z_j)^{N_{ij}}},\qquad
 P(z_2+t,z_2,z_3,\ldots)=\sum_{\ell\ge0}P_\ell(z_2,z_3,\ldots)t^\ell.
\]
The second sum is finite.  For \(B>0\), formal inversion gives
\begin{equation}\label{op:eq:binomial-pole}
 (x+t)^{-B}
 =\sum_{h\ge0}(-1)^h\binom{B+h-1}{h}x^{-B-h}t^h.
\end{equation}
For \(B=0\) the series is \(1\), with only its constant term.
After substitution the sole negative power of \(t\) outside these
power series is \(t^{-A}\).  A contribution to \(t^q\) therefore
satisfies
\(\ell+\sum h_j=q+A=r\).
There are finitely many such tuples.  Their denominators are exactly
those in \eqref{op:eq:collision-budget}.  This proves both the vanishing
and the upper bounds.

Fix \(j\) with \(B_j>0\), and take \(P=1\).  The term with
\(h_j=r\) and every other \(h_i=0\) has pole order
\(B_j+C_j+r\) along \(x_j=0\).  Its coefficient
\((-1)^r\binom{B_j+r-1}{r}\) is nonzero in characteristic zero.
Every other term has smaller order on that diagonal.  At its generic
point all the remaining differences are units, so cancellation of
this leading pole is impossible.  If \(B_j=0\), use instead
\(P=(z_1-z_2)^r\).  Only the constant terms of all the binomial
series contribute, giving exact order \(C_j\).  This proves the
minimality assertion, including zero exponents.
\end{proof}

For three points the coefficient realizing the positive case is
\begin{equation}\label{op:eq:sharp-three-point}
 [t^{r-A}]\,t^{-A}(x+t)^{-B}x^{-C}
 =(-1)^r\binom{B+r-1}{r}x^{-(B+C+r)},
 \qquad B>0.
\end{equation}
When \(A\ge1\), a residue uses \(r=A-1\).  When \(A=0\),
the residue is zero.  In the constant-term operation \(q=0\), the
extra budget is \(r=A\).  Formula~\eqref{op:eq:collision-budget} also
shows that different remaining diagonals share this budget.
Replacing it by independent maxima loses that information.

The proof of the upper bound uses integer binomial coefficients and
works over every field.  Sharpness can fail in positive characteristic.
For example, if \(\operatorname{char}k=p>0\), then
\[
 (x+t)^{-p}=(x^p+t^p)^{-1},
\]
so its coefficients at \(t^r\), \(1\le r<p\), vanish.

\subsection{Exact coefficient support in positive characteristic}
\label{op:subsec:positive-characteristic-poles}

In this subsection only, \(k\) may be a field of arbitrary
characteristic.  Fix \(n\ge3\) and algebraically independent
variables \(x_3,\ldots,x_n\) over \(k\).  Put
\(K=k(x_3,\ldots,x_n)\).  For nonnegative integers
\(a,b_j,c_j\), consider the monomial rational model
\begin{equation}\label{op:eq:monomial-collision-model}
 f(t)=t^{-a}\prod_{j=3}^n(x_j+t)^{-b_j}x_j^{-c_j}
 \in K((t)),\qquad d_j=b_j+c_j,\quad D=\sum_{j=3}^nd_j.
\end{equation}
Every \(x_j+t\) is inverted in \(K[[t]]\) before the factor
\(t^{-a}\) is applied.  No completion in the \(x_j\) variables
is used.  This model has polynomial numerator one and no
denominator between two remaining labels.  Write
\(\mathcal C_m=[t^{-a+m}]f(t)\) for \(m\ge0\).

Define integers
\[
 H_b(r)=
 \begin{cases}
 \binom{b+r-1}{r},&b>0,\\
 1,&b=0,\ r=0,\\
 0,&b=0,\ r>0.
 \end{cases}
\]
Their use as coefficients means reduction under \(\mathbb Z\to k\).
This definition requires no factorial denominator in \(k\).
Let
\[
 T_b(k)=\{r\ge0:H_b(r)\ne0\text{ in }k\}.
\]

\begin{lemma}[The digit criterion]
\label{op:lem:pole-no-carry}
Suppose \(\operatorname{char}k=p>0\) and \(b>0\).
Then \(r\in T_b(k)\) if and only if the addition of \(r\)
and \(b-1\) in base \(p\) has no carry.  In all
characteristics, \(T_0(k)=\{0\}\).  In characteristic zero,
\(T_b(k)=\mathbb N\) for \(b>0\).
\end{lemma}

\begin{proof}
For \(0<i<p\), the integer \(\binom pi\) is divisible by
\(p\).  Hence \((1+z)^p=1+z^p\) over \(k\).
If \(N=\sum_\ell N_\ell p^\ell\) is the base-\(p\)
expansion of a nonnegative integer, repeated Frobenius gives
\[
 (1+z)^N=\prod_\ell(1+z^{p^\ell})^{N_\ell}.
\]
Every exponent chosen from the \(\ell\)-th factor is
\(s_\ell p^\ell\), with \(0\le s_\ell\le N_\ell<p\).
These choices have a unique base-\(p\) sum.  Thus, for
\(r=\sum_\ell r_\ell p^\ell\), the coefficient of \(z^r\)
is zero unless \(r_\ell\le N_\ell\) for all \(\ell\).
When these inequalities hold, the coefficient is
\(\prod_\ell\binom{N_\ell}{r_\ell}\), which is nonzero.
Indeed, the factorial arguments in these digit binomials are less
than \(p\), so all their factors are invertible in \(k\).

Apply this result to \(N=r+b-1\).  If the addition has no
carry, then \(N_\ell=r_\ell+(b-1)_\ell\ge r_\ell\)
in every digit.  If a carry occurs, let \(\ell\) be its first
position.  There is no incoming carry there, so
\[
 N_\ell=r_\ell+(b-1)_\ell-p<r_\ell.
\]
The corresponding binomial coefficient vanishes.  This proves the
criterion for \(H_b(r)=\binom{r+b-1}{r}\).
The assertions for \(b=0\) and characteristic zero follow from
the defining integers.
\end{proof}

\begin{theorem}[Exact support and divisor orders]
\label{op:thm:monomial-pole-support}
For the model \eqref{op:eq:monomial-collision-model}, all coefficients
below \(t^{-a}\) vanish, and
\begin{equation}\label{op:eq:exact-monomial-coefficient}
 \mathcal C_m=(-1)^m
 \sum_{\substack{r_3,\ldots,r_n\ge0\\\sum_jr_j=m}}
 \left(\prod_{j=3}^nH_{b_j}(r_j)\right)
 \prod_{j=3}^nx_j^{-(d_j+r_j)}.
\end{equation}
Define the finite set
\[
 \mathcal R_m=
 \{(r_3,\ldots,r_n):\textstyle\sum_jr_j=m,
                    \ r_j\in T_{b_j}(k)\}.
\]
Then
\begin{equation}\label{op:eq:exact-coefficient-support}
 \mathcal C_m\ne0
 \quad\Longleftrightarrow\quad
 \mathcal R_m\ne\varnothing
 \quad\Longleftrightarrow\quad
 m\in T_{b_3}(k)+\cdots+T_{b_n}(k).
\end{equation}
The last expression is a sum of sets.
For a nonzero coefficient, its pole order \(p_j\) along
\(x_j=0\) is exactly
\begin{equation}\label{op:eq:exact-divisorial-pole}
 p_j(\mathcal C_m)=d_j+\max_{\mathbf r\in\mathcal R_m}r_j.
\end{equation}
If \(b_j=0\), this order is \(c_j\).
If \(b_j>0\), the bound \(d_j+m\) is attained if and
only if \(m\in T_{b_j}(k)\).  In characteristic \(p\),
this is precisely the no-carry criterion of
Lemma~\ref{op:lem:pole-no-carry} with \(r=m\).

Finally, every nonzero \(\mathcal C_m\) is homogeneous of degree
\(D+m\) in the polynomial variables \(\xi_j=x_j^{-1}\).
This common degree of the surviving monomials is distinct from
the sum of their separate divisor pole maxima.
\end{theorem}

\begin{proof}
The geometric series for \((1+z)^{-1}\), multiplied \(b\)
times, gives
\[
 (1+z)^{-b}=\sum_{r\ge0}(-1)^rH_b(r)z^r.
\]
For \(b>0\), the coefficient counts weak compositions of
\(r\) into \(b\) parts.  The same formula for \(b=0\)
is the constant series one.  This proves the identity first over
the integers and then over \(k\).  Multiplying the finitely many
factors in \eqref{op:eq:monomial-collision-model} proves
\eqref{op:eq:exact-monomial-coefficient}.  At fixed \(m\), only
finitely many tuples occur.

Different tuples give different Laurent monomials in the independent
variables \(x_j\).  A tuple survives exactly when every scalar
factor \(H_{b_j}(r_j)\) is nonzero.  Distinct surviving
monomials are linearly independent, proving
\eqref{op:eq:exact-coefficient-support}.

Fix \(j\) and put \(R=\max_{\mathbf r\in\mathcal R_m}r_j\).
View the coefficient as a Laurent polynomial in \(x_j\) over
\(k(x_i:i\ne j)\).  The terms with \(r_j=R\) have lowest
exponent \(-d_j-R\).  Their coefficient is a nonzero Laurent
polynomial in the remaining independent variables.  Thus they
cannot cancel, proving \eqref{op:eq:exact-divisorial-pole}.
When \(b_j=0\), only \(r_j=0\) is allowed.
When \(b_j>0\), equality \(r_j=m\) forces all other
\(r_i\) to be zero.  This tuple survives exactly when
\(H_{b_j}(m)\ne0\), since \(H_{b_i}(0)=1\) for every
\(i\).  This proves the saturation criterion.

Each surviving monomial has degree
\(\sum_j(d_j+r_j)=D+m\) in the \(\xi_j\).
Their sum is nonzero and homogeneous of that degree.
\end{proof}

In characteristic zero, if \(s\) of the \(b_j\) are positive,
then every nonzero coefficient satisfies
\begin{equation}\label{op:eq:sum-divisor-maxima}
 \sum_jp_j(\mathcal C_m)=D+s\,m.
\end{equation}
For \(s=0\), only \(m=0\) occurs.
In characteristic \(p\), a positive \(b_j\) still gives
infinitely many nonzero coefficients and unbounded pole orders.
Choose \(P=p^L>b_j-1\).  Every multiple \(m\) of \(P\)
adds to \(b_j-1\) without a carry.  The tuple \(r_j=m\),
\(r_i=0\) for \(i\ne j\), therefore survives and attains
pole order \(d_j+m\).

\begin{example}[Four points and the two degree counts]
\label{op:ex:four-point-degree-counts}
Let \(n=4\), \(x=x_3\), \(y=x_4\),
\(b_3=b_4=1\), and \(c_3=c_4=0\).  Over every field,
\[
 \mathcal C_1=-x^{-2}y^{-1}-x^{-1}y^{-2}
             =-\frac{x+y}{x^2y^2}.
\]
Its pole orders are two along each coordinate divisor.
Their sum is four, although both Laurent monomials have degree
three in \(x^{-1},y^{-1}\).  Thus \(D+m\) does not bound
the sum of the separate divisor maxima.
Algebraic independence is essential: after setting \(y=x\),
this coefficient becomes \(-2x^{-3}\) and vanishes in
characteristic two.
\end{example}

\begin{example}[Four points in characteristic two]
\label{op:ex:four-point-characteristic-two}
Let \(n=4\), \(b_3=b_4=3\), and \(c_3=c_4=0\).
In characteristic two, the digit criterion gives
\(T_3(k)=\{r\ge0:r\equiv0,1\pmod4\}\), since \(3-1\)
has binary digits \(10\).  Equivalently,
\[
 (1+z)^{-3}=\frac{1+z}{1+z^4}
           =\sum_{h\ge0}(z^{4h}+z^{4h+1}).
\]
With \(x=x_3\) and \(y=x_4\), the coefficient formula gives
\[
 \begin{aligned}
 \mathcal C_1&=x^{-4}y^{-3}+x^{-3}y^{-4},\\
 \mathcal C_2&=x^{-4}y^{-4},\\
 \mathcal C_3&=0.
 \end{aligned}
\]
At \(m=2\), the only surviving allocation is \((1,1)\).
Both pole orders drop from the bound five to four, while the
degree \(D+m=8\) remains exact.  At \(m=3\), no allocation
survives.  The set \(T_3(k)\) is not closed under addition:
it contains one but not two.  The set sums in
\eqref{op:eq:exact-coefficient-support} require no such closure.
\end{example}

\begin{example}[A pole drops while every coefficient survives]
\label{op:ex:four-point-prime-staircase}
Let \(\operatorname{char}k=p>0\), \(n=4\),
\(b_3=p\), \(b_4=1\), and \(c_3=c_4=0\).
Here \(T_p(k)=p\mathbb N\) and \(T_1(k)=\mathbb N\).
Every \(m\ge0\) is supported.  Among the surviving pairs,
the largest first entry is \(p\lfloor m/p\rfloor\),
and the largest second entry is \(m\).  Hence
\[
 p_3(\mathcal C_m)=p+p\left\lfloor\frac mp\right\rfloor,
 \qquad p_4(\mathcal C_m)=1+m.
\]
The first divisor attains its bound \(p+m\) exactly when
\(p\mid m\).  Failure of one saturation criterion therefore
need not make the whole coefficient vanish.
\end{example}

These exact support statements concern the monomial input
\eqref{op:eq:monomial-collision-model}.  Polynomial numerators, sums
of rational inputs, and variable specializations require their own
cancellation checks.  The characteristic-zero hypotheses on the
simplex-form and lattice constructions below remain in force.

\subsection{Order faces and the completion domain}

Let \(\operatorname{Ord}(I)\) be the finite set of total orders of
\(I\).  For a nonempty subset \(\Sigma\), let \(R_\Sigma\)
invert exactly the differences whose comparisons vary among the
orders in \(\Sigma\).  Define \(C_I\) to consist of polynomial
forms on the simplex with vertices \(\operatorname{Ord}(I)\), with
coefficients in \(L_I\), whose restriction to each face \(\Sigma\)
has coefficients in \(R_\Sigma\).  The differential acts only on
simplex coordinates.

\begin{lemma}[Pair faces detect permitted poles]
\label{op:lem:pair-face-poles}
A polynomial simplex form belongs to \(C_I\) if and only if its
restriction to the face of orders putting \(i\) before \(j\) has
no pole along \(z_i=z_j\), for every ordered pair \(i,j\).
\end{lemma}

\begin{proof}
Necessity is part of the definition.  For sufficiency, suppose all
orders in \(\Sigma\) put \(i\) before \(j\).  This face lies
inside the specified pair face, so restriction excludes that pole.
Write each coefficient in reduced form in the polynomial fraction
field.  The nonassociated prime differences in its denominator are
precisely its possible diagonal poles.  Excluding each forbidden
difference therefore excludes them simultaneously, giving
coefficients in \(R_\Sigma\).
\end{proof}

For a consecutive block \(\{1,2\}\), the substitution
\(z_1=z_2+t\) is available on the associated order overlaps.
An order \(\sigma\) of the merged labels and an inner order \(\tau\)
give a leaf order \(\operatorname{lex}(\sigma,\tau)\) by replacing
the merged label with the consecutive ordered pair.  This map is
injective.  If \(s_\rho\), \(v_\sigma\), and \(w_\tau\)
are the respective barycentric coordinates, the simplex pullback is
\[
 s_{\operatorname{lex}(\sigma,\tau)}\longmapsto v_\sigma w_\tau,
 \qquad s_\rho\longmapsto0
 \quad\text{for nonconsecutive leaf orders}.
\]
Send differentials to differentials of these expressions.
The images sum to \((\sum v_\sigma)(\sum w_\tau)=1\), so this
defines a map of polynomial differential-form algebras.
An inverted internal difference requires the two inner orders and
maps to \(t\).  An inverted cross difference requires varying outer
orders, so its constant term \(x_j\) is already a unit.
Thus Theorem~\ref{op:thm:collision-poles} applies coefficientwise on
each such overlap.  Pullback of polynomial simplex forms adds no
spatial denominator.  On a face that fixes an outer comparison,
Lemma~\ref{op:lem:pair-face-poles} removes the corresponding cross poles
before substitution.  This is a bound for the permitted coefficient
maps on those faces.

The order of localization and completion matters even for one cross
difference.  There is a strict inclusion
\begin{equation}\label{op:eq:pole-completion-inclusion}
 k[x][[t]][x^{-1}]
 \;\subsetneq\;
 k[x,x^{-1}][[t]].
\end{equation}
Every series on the left has a common lower bound on its powers of
\(x\).  The inverse
\[
 (x+t)^{-1}=\sum_{h\ge0}(-1)^h x^{-h-1}t^h
\]
belongs to the right side and has no such bound.  Injectivity and
strictness follow by comparing coefficients of \(t\).
Consequently, a fixed rational function can have collision
coefficients with unbounded outer pole orders.

Repeated binary coefficient extractions along a specified ordered
tree remain rational functions.  Applying
Theorem~\ref{op:thm:collision-poles} successively gives finite bounds
at every vertex and on every final coefficient.  Full scale series
require a common domain in addition to these finite bounds.
For instance
\[
 \sum_{h\ge0}\epsilon^h\eta^{-h-1}
 \notin k[[\epsilon,\eta]][(\epsilon\eta)^{-1}],
\]
because its \(\eta\)-exponents have no lower bound.  Passing through
arbitrary intermediate Laurent completions is therefore not licensed
by the coefficient theorem.  Formal unit specialization also needs
care: no unital ring homomorphism from \(k[[\epsilon]]\) to \(k\)
can send \(\epsilon\) to \(1\), since \(1-\epsilon\) is a unit.

\subsection{Sharpness for actual vertex products}

For the vertex-algebra realization in this subsection, assume
\(\operatorname{char}k=0\).
The same exponent occurs in vertex locality.  Use the mode convention
\[
 Y(a,z)=\sum_{n\in\mathbb Z}a_{(n)}z^{-n-1}.
\]
Let \(N(a,b)\) be the least nonnegative integer annihilating
\([Y(a,z),Y(b,w)]\) after multiplication by \((z-w)^{N(a,b)}\).
All states in the example below are even.  The product \(a_{(n)}b\)
is the coefficient \([t^{-n-1}]Y(a,t)b\).  Roitman's quantitative
Dong lemma gives, for \(n<A\),
\[
 N(a_{(n)}b,c)\le
 \begin{cases}
 B+C+A-n-1,&B>0,\\
 C,&B=0,
 \end{cases}
\]
when \(N(a,b)\le A\), \(N(a,c)\le B\), and
\(N(b,c)\le C\).  See \cite{Roi02}*{Section 7, Lemma 1,
equation (16), pp.\ 14--15 of the preprint}.

Here is a realization of equality for every
\(A,C\ge0\), \(B>0\), and \(r=A-n-1\ge0\).
Let \(\Lambda\) have basis \(\alpha,\beta,\gamma\), with
\[
 (\alpha,\beta)=-A,\quad (\alpha,\gamma)=-B,\quad
 (\beta,\gamma)=-C.
\]
Give each basis vector squared length \(2(A+B+C+1)\).
The integral Gram matrix is strictly diagonally dominant and positive
definite, with even diagonal.  Thus \(\Lambda\) is an even lattice.
Its vertex algebra has underlying space
\[
 V_\Lambda=
 \Sym\bigl(\mathfrak h\otimes z^{-1}k[z^{-1}]\bigr)
 \otimes k_\varepsilon[\Lambda],\qquad
 \mathfrak h=k\otimes_{\mathbb Z}\Lambda.
\]
The group algebra is twisted by a bimultiplicative sign
\(\varepsilon\) satisfying
\(\varepsilon(\lambda,\mu)/\varepsilon(\mu,\lambda)
=(-1)^{(\lambda,\mu)}\).
The Heisenberg modes satisfy
\([h(m),h'(n)]=m(h,h')\delta_{m,-n}\).
The vacuum in charge \(\lambda\) is denoted by \(e^\lambda\).
With positive modes annihilating this vacuum, the lattice field is
\[
 Y(e^\lambda,z)=
 \exp\left(\sum_{m>0}\frac{\lambda(-m)}m z^m\right)
 \exp\left(-\sum_{m>0}\frac{\lambda(m)}m z^{-m}\right)
 e^\lambda z^{\lambda(0)}.
\]
The lattice construction, vacuum products, and charge grading are
described in \cite{Roi02}*{Section 4, pp.\ 8--10 of the preprint}.
In particular,
\[
 Y(e^\alpha,t)e^\beta
 =\varepsilon(\alpha,\beta)t^{-A}
 \exp\left(\sum_{m>0}\frac{\alpha(-m)}m t^m\right)
 e^{\alpha+\beta}.
\]

Let \(p_\lambda\) take the coefficient of \(e^\lambda\) and
kill all positive oscillator levels and other charges.  Commuting
positive modes through negative modes in the field formula gives
\[
 p_{\alpha+\beta+\gamma}
\bigl(Y(e^\alpha,z_1)Y(e^\beta,z_2)e^\gamma\bigr)
 =\epsilon_0\,\iota_{12}
 \bigl((z_1-z_2)^{-A}z_1^{-B}z_2^{-C}\bigr),
\]
where \(\iota_{12}\) expands in nonnegative powers of \(z_2/z_1\),
and \(\epsilon_0=\varepsilon(\beta,\gamma)
\varepsilon(\alpha,\beta+\gamma)\in\{1,-1\}\).
Indeed, the commutator of the two exponential exponents is
the scalar series producing \((z_1-z_2)^{-A}\).
The zero modes supply \(z_1^{-B}z_2^{-C}\), and the remaining
negative-mode exponentials project to \(1\).
In the rational function set \(z_1=x+t\), \(z_2=x\), then
expand at \(t=0\) with \(x\) invertible.
Vertex associativity identifies this expansion with the iterate
\(Y(Y(e^\alpha,t)e^\beta,x)e^\gamma\).
Formula~\eqref{op:eq:sharp-three-point} consequently gives
\begin{equation}\label{op:eq:lattice-pole-witness}
 p_{\alpha+\beta+\gamma}
 \bigl(Y((e^\alpha)_{(A-r-1)}e^\beta,x)e^\gamma\bigr)
 =\epsilon_0(-1)^r\binom{B+r-1}{r}x^{-B-C-r}.
\end{equation}

For completeness, the full state-valued product has no higher pole.
A charge-\(\lambda\) vector has conformal weight
\((\lambda,\lambda)/2+\ell\), with oscillator level \(\ell\ge0\).
The merged state has charge \(\alpha+\beta\) and level \(r\).
A term \(x^{-m-1}\) acting on \(e^\gamma\) therefore has level
\(B+C+r-m-1\) in charge \(\alpha+\beta+\gamma\).
Negative levels are absent.  Hence its pole order is at most
\(B+C+r\), while \eqref{op:eq:lattice-pole-witness} attains that
order.  In a vertex algebra this maximal nonzero nonnegative product
index plus one is the locality order, by the field commutator formula.
Thus the bound is attained by actual fields, including the residue
example \(A=B=C=1\), \(r=0\).

These coefficient and field calculations specify the filtration
needed by a collision presentation.  They supply neither the operadic
lift in Proposition~\ref{prop:locality-lift} nor the monoidal functor
in Problem~\ref{prob:carrier-change}.  The Euler product below uses
root weights and homological parity, independently of this pole
filtration.
